\section{附录：核心算法代码摘录}
\label{app:core-code}

下列代码摘自随文提交的计算脚本，仅展示关键算法；数据摄入、绘图和结果写出等实现不在纸质附录展开。完整程序见电子材料的 \texttt{code/}，问题一的语义解码和固定参照评分见 \texttt{quality\_scoring.py} 与 \texttt{q1\_pipeline.py}。

% 以下 firstline/lastline 对应随文提交的源文件；修改源文件时须同步核对范围。

\subsection{问题一：质量评分与配比回归}

\lstinputlisting[style=pythonstyle,firstline=41,lastline=75,caption={二分类与六分类信号的语义解码},label={code:q1-score}]{../code/quality_scoring.py}

\lstinputlisting[style=pythonstyle,firstline=117,lastline=142,caption={A1 固定百分位与三组平衡评分},label={code:q1-quality}]{../code/quality_scoring.py}

\lstinputlisting[style=pythonstyle,firstline=54,lastline=59,caption={高教育价值与广告判定的冲突条件},label={code:q1-conflict}]{../code/q1_pipeline.py}

\lstinputlisting[style=pythonstyle,firstline=55,lastline=72,caption={单纯形配比回归及可识别系数编码},label={code:q1-mixture}]{../code/problem1.py}

\subsection{问题二：经典与广义标度律}

\lstinputlisting[style=pythonstyle,firstline=37,lastline=67,caption={经典标度律的对数域稳健拟合},label={code:q2-classic}]{../code/problem2.py}

\lstinputlisting[style=pythonstyle,firstline=70,lastline=117,caption={广义标度律及质量指数拟合},label={code:q2-generalized}]{../code/problem2.py}

\subsection{问题三：预算约束下的资源配置}

\lstinputlisting[style=pythonstyle,firstline=114,lastline=169,caption={预算消元、质量剖面搜索与边界检验},label={code:q3-solver}]{../code/problem3.py}

\subsection{问题四：能力前沿与情景预测}

\lstinputlisting[style=pythonstyle,firstline=190,lastline=211,caption={单调有界的损失到能力桥接},label={code:q4-frontier}]{../code/problem4.py}

\lstinputlisting[style=pythonstyle,firstline=325,lastline=351,caption={联合重拟合截距与斜率的自助情景预测},label={code:q4-forecast}]{../code/problem4.py}
